%% kapa: exam=1 round=35 year=2024 period=1 form=A subject=economics subject_ko=경제학원론 seq=2 number=42
%% source: Q-Net 감정평가사 기출문제 (cst00309) artlSeq=5213229 file=raw/1008/5213229/1교시 문제지 원본.pdf
%% conversion: kordoc 4.24.51 -> tools/gen1.py (auto); review=proofread-sonnet; flags=none
\begin{question}{42}
재화의 특성에 관한 설명으로 옳은 것은?
\begin{choices}
  \item 사치재는 수요의 가격탄력성이 1보다 큰 재화를 말한다.
  \item 열등재는 가격이 오르면 수요가 감소하는 재화를 말한다.
  \item 절댓값으로 볼 때, 가격효과가 소득효과보다 큰 열등재를 기펜재(Giffen goods)라고 한다.
  \item 두 상품이 완전 대체재이면 무차별 곡선은 원점에 대하여 볼록한 모양이다.
  \item 수요가 가격 탄력적인 상품을 판매하는 기업이 가격을 내리면 판매수입은 증가한다.
\end{choices}
\end{question}
\answer{5}
